\documentclass[12pt,a4paper]{article}

% ------------------------------------------------------------------
% Pacotes
% ------------------------------------------------------------------
\usepackage[utf8]{inputenc}
\usepackage[T1]{fontenc}
\usepackage[brazil]{babel}
\usepackage[a4paper,left=3cm,right=2cm,top=3cm,bottom=2cm]{geometry}
\usepackage{setspace}
\usepackage{indentfirst}
\usepackage{titlesec}
\usepackage{enumitem}
\usepackage{graphicx}
\usepackage{hyperref}

\hypersetup{
  colorlinks=true,
  linkcolor=black,
  urlcolor=blue,
  citecolor=black
}

% Espaçamento 1,5 conforme normas acadêmicas
\onehalfspacing

% Formatação dos títulos de seção
\titleformat{\section}{\normalfont\large\bfseries}{\thesection}{1em}{}
\titleformat{\subsection}{\normalfont\normalsize\bfseries}{\thesubsection}{1em}{}

% ------------------------------------------------------------------
% Capa
% ------------------------------------------------------------------
\begin{document}

\begin{titlepage}
  \centering
  {\large UNIVERSIDADE PRESBITERIANA MACKENZIE\par}
  \vspace{0.5cm}
  {\normalsize Administração de Negócios\par}
  \vfill
  {\Large\bfseries Diagnóstico Organizacional da Sentinela Cyber Consultoria em Segurança da Informação\par}
  \vspace{0.4cm}
  {\large Relatório dissertativo aplicando conceitos de Administração a uma empresa de consultoria em segurança da informação\par}
  \vfill
  {\normalsize\bfseries Autores:\par}
  {\normalsize Bruno Viana\par}
  {\normalsize Felipe Koike\par}
  {\normalsize Jônatas Brito\par}
  \vfill
  {\normalsize Avaliação N2\par}
  {\normalsize São Paulo, 2026\par}
\end{titlepage}

% ------------------------------------------------------------------
% Sumário
% ------------------------------------------------------------------
\tableofcontents
\newpage

% ==================================================================
% INTRODUÇÃO
% ==================================================================
\section{Introdução}

A transformação digital ampliou de forma acelerada a superfície de exposição
das organizações a ameaças cibernéticas. Nesse cenário, ataques de engenharia
social, em especial o \textit{phishing}, tornaram-se uma das principais portas
de entrada para incidentes de segurança, explorando não apenas falhas
tecnológicas, mas sobretudo o comportamento humano. É justamente nessa
fronteira entre tecnologia e pessoas que se posiciona a empresa analisada neste
relatório.

O presente trabalho tem por objetivo aplicar os conceitos estudados na
disciplina de Administração de Negócios à realidade de uma empresa de
consultoria em segurança da informação, a Sentinela Cyber Consultoria. A
análise percorre a apresentação da organização, o diagnóstico de seu ambiente
interno, a forma como executa o processo administrativo, o papel das
\textit{soft skills} e do relacionamento interpessoal, a relação entre suas
práticas e a evolução do pensamento administrativo e, por fim, uma reflexão
crítica que confronta teoria e prática, propondo recomendações de gestão.

A metodologia adotada combina a revisão dos conceitos teóricos vistos em sala,
referentes às escolas Clássica, das Relações Humanas e Sistêmica, com a
caracterização de uma empresa de pequeno porte do setor de tecnologia,
permitindo observar como princípios administrativos consagrados continuam (ou
não) aplicáveis a um negócio jovem, intensivo em conhecimento e fortemente
dependente de capital humano qualificado.

% ==================================================================
% 1. APRESENTAÇÃO DA EMPRESA
% ==================================================================
\section{Apresentação da Empresa Pesquisada}

\subsection{Breve histórico e área de atuação}

A Sentinela Cyber Consultoria é uma empresa de tecnologia que atua no mercado
de segurança da informação na modalidade de consultoria. Foi fundada por três
alunos da Universidade Presbiteriana Mackenzie (Bruno Viana, Felipe Koike e
Jônatas Brito), que uniram suas trajetórias profissionais complementares para
oferecer ao mercado serviços voltados à proteção do elo mais frágil da cadeia
de segurança: as pessoas.

A complementaridade dos fundadores é um dos pilares do negócio. Bruno Viana
traz dez anos de experiência na área de auditoria, fundamentais para a
estruturação de processos de \textit{compliance} e avaliação de riscos. Felipe
Koike soma nove anos como analista de dados sênior em uma instituição
bancária, competência decisiva para a análise de indicadores e métricas de
segurança. Jônatas Brito reúne quatro anos em infraestrutura de TI e três anos
como desenvolvedor \textit{backend}, garantindo a base técnica para a
condução de testes, simulações e desenvolvimento de soluções internas.

A empresa atua de forma especializada na prestação dos seguintes serviços:

\begin{itemize}[noitemsep]
  \item Treinamentos antiphishing e campanhas de conscientização (\textit{security awareness});
  \item Capacitação em boas práticas no uso de ferramentas da internet e de produtividade corporativa;
  \item Consultoria em políticas de segurança da informação e \textit{compliance};
  \item Simulações de ataques de engenharia social e geração de relatórios de maturidade em segurança.
\end{itemize}

O posicionamento da empresa é, portanto, o de uma consultoria \textit{boutique}
de segurança da informação, com forte ênfase no fator humano, atendendo
organizações que buscam reduzir incidentes decorrentes de erro humano e
fortalecer sua cultura de segurança.

\subsection{Estrutura organizacional básica}

A Sentinela Cyber organiza-se em três áreas, cada uma liderada por um dos
sócios fundadores e estruturada de forma padronizada com sete integrantes: o
diretor da área, um gerente de projetos (PM) ou \textit{product owner} (PO),
quatro colaboradores e um estagiário. Esse desenho equilibra as três frentes,
totaliza vinte e uma pessoas e cria um caminho claro de desenvolvimento de
carreira dentro de cada equipe. A divisão de responsabilidades acompanha as
especialidades de cada fundador:

\begin{itemize}[noitemsep]
  \item Diretoria de Auditoria e Compliance (Bruno Viana): responsável pela
        metodologia de avaliação de riscos, conformidade e relacionamento com
        áreas de governança dos clientes;
  \item Diretoria de Dados e Inteligência (Felipe Koike): responsável pela
        análise de indicadores, mensuração de resultados das campanhas e apoio à
        tomada de decisão baseada em dados;
  \item Diretoria de Tecnologia e Operações (Jônatas Brito): responsável pela
        infraestrutura, pelas plataformas de simulação de phishing e pelo
        desenvolvimento das ferramentas internas.
\end{itemize}

Dentro de cada área, os papéis cumprem funções complementares. O diretor
responde pela estratégia da frente e pelo relacionamento com os clientes. O PM
ou PO coordena o fluxo de trabalho, prioriza as entregas e faz a interface
entre a equipe técnica e os clientes. Os quatro colaboradores executam as
atividades-fim da área, como auditorias, análises de dados ou operações
técnicas. O estagiário apoia as tarefas operacionais e está em formação,
sustentando o pilar de Evolução Contínua e a renovação de talentos da empresa.

Trata-se de uma estrutura predominantemente funcional, organizada por
especialidade, na qual a padronização do tamanho e da composição das equipes
favorece a coordenação, a previsibilidade da capacidade produtiva e a
escalabilidade do negócio.

% ==================================================================
% 2. DIAGNÓSTICO DO AMBIENTE ORGANIZACIONAL
% ==================================================================
\section{Diagnóstico do Ambiente Organizacional}

\subsection{Cultura organizacional: valores, clima e práticas de liderança}

A cultura da Sentinela Cyber é fortemente influenciada pela origem técnica e
acadêmica de seus fundadores e pela própria natureza do negócio, que exige
disciplina, ética e confiança. Para traduzir esses elementos, a empresa
estrutura sua cultura em quatro pilares que expressam responsabilidade
conjunta, união entre colaboradores e boas práticas no ambiente empresarial:

\begin{enumerate}[label=\arabic*.]
  \item Defesa Coletiva: a segurança é responsabilidade de todos; cada
        colaborador é corresponsável pela proteção da empresa e dos clientes,
        reforçando o princípio da responsabilidade conjunta.
  \item Confiança em Rede: a união entre os colaboradores é tratada como um
        ativo estratégico; a colaboração aberta, o compartilhamento de
        conhecimento e o apoio mútuo sustentam os resultados coletivos.
  \item Integridade Inegociável: a ética, a transparência e o sigilo são
        tratados como condição inegociável, traduzindo as boas práticas no
        ambiente empresarial e na relação com os clientes.
  \item Evolução Contínua: diante de um cenário de ameaças em constante mutação,
        o aprendizado permanente e a melhoria contínua são valores centrais para
        a sustentabilidade do negócio.
\end{enumerate}

O clima organizacional, favorecido pela proximidade entre as equipes, é de
abertura, autonomia e comunicação direta. A liderança é exercida de forma
participativa: as decisões estratégicas são debatidas coletivamente entre os
três sócios, e cada um lidera tecnicamente sua área de especialidade. Esse
modelo aproxima-se de um estilo de liderança democrática, no qual a autoridade
decorre mais da competência reconhecida do que da posição hierárquica.

\subsection{Principais desafios enfrentados pela empresa}

Apesar do diferencial competitivo, a Sentinela Cyber enfrenta desafios típicos
de uma empresa jovem e de um setor altamente dinâmico:

\begin{itemize}[noitemsep]
  \item Dependência das pessoas-chave: por concentrarem a estratégia e o
        conhecimento técnico das áreas, os três sócios fundadores ainda
        representam um ponto crítico para a continuidade do negócio;
  \item Construção de credibilidade: por ser nova no mercado, a empresa precisa
        conquistar a confiança de clientes em um setor onde reputação é
        determinante;
  \item Atualização constante: a evolução acelerada das ameaças exige
        investimento permanente em capacitação e atualização técnica;
  \item Formalização de processos: a informalidade característica do início do
        negócio precisa ceder lugar a processos documentados e escaláveis;
  \item Mensuração de resultados intangíveis: demonstrar o retorno de
        treinamentos comportamentais é um desafio metodológico recorrente.
\end{itemize}

% ==================================================================
% 3. PROCESSO ADMINISTRATIVO NA PRÁTICA
% ==================================================================
\section{Processo Administrativo na Prática}

O processo administrativo, conforme proposto pela abordagem clássica e
consolidado na literatura, compreende quatro funções interligadas:
planejamento, organização, direção e controle. A seguir, descreve-se como cada
uma se manifesta na Sentinela Cyber.

\subsection{Planejamento}

O planejamento estratégico da empresa parte da definição de seu nicho, a
consultoria em segurança da informação centrada no fator humano. No nível
tático, são estabelecidas metas trimestrais de aquisição de clientes e de
desenvolvimento de novos pacotes de serviços. No nível operacional, cada
projeto de consultoria conta com um plano de ação com escopo, cronograma e
entregáveis definidos junto ao cliente.

\subsection{Organização}

A organização se materializa na divisão de responsabilidades entre as três
diretorias (Auditoria/Compliance, Dados/Inteligência e Tecnologia/Operações),
na padronização de metodologias de trabalho e na alocação de recursos
tecnológicos. Cada projeto é organizado em etapas: diagnóstico inicial,
execução das campanhas e simulações, análise de resultados e plano de melhoria.

\subsection{Direção}

A direção é exercida por meio da liderança participativa dos sócios, da
coordenação direta das equipes de projeto e da manutenção de um ambiente de
motivação e engajamento. A comunicação fluida e a autonomia concedida aos
colaboradores são instrumentos centrais de direção em uma organização
intensiva em conhecimento.

\subsection{Controle}

O controle ocorre por meio do acompanhamento de indicadores de desempenho dos
clientes, como a taxa de cliques em campanhas de phishing simulado, a evolução
da taxa de reporte de e-mails suspeitos e o índice de conclusão dos
treinamentos. Internamente, são monitorados prazos, satisfação dos clientes e
resultados financeiros, permitindo ajustes nas ações planejadas.

\subsection{Ferramentas de apoio à decisão}

A empresa utiliza um conjunto de ferramentas que sustentam a tomada de decisão
baseada em evidências:

\begin{itemize}[noitemsep]
  \item Painéis (\textit{dashboards}) de indicadores, sob responsabilidade da
        diretoria de dados, que consolidam métricas de segurança e desempenho;
  \item Plataformas de simulação de phishing, que geram dados objetivos sobre o
        comportamento dos usuários;
  \item Sistemas de gestão de projetos (ferramentas \textit{kanban}), para
        acompanhamento de tarefas e prazos;
  \item Matrizes de risco e análise SWOT, empregadas no diagnóstico estratégico
        de clientes e da própria empresa.
\end{itemize}

% ==================================================================
% 4. RELACIONAMENTO INTERPESSOAL E SOFT SKILLS
% ==================================================================
\section{Relacionamento Interpessoal e Soft Skills}

\subsection{Motivação, comunicação e trabalho em equipe}

Em uma consultoria formada por sócios com forte vínculo acadêmico e
profissional, a motivação é sustentada pelo sentimento de propósito comum e
pela autonomia técnica de cada um. A comunicação é direta e horizontal,
favorecida pela proximidade entre as equipes, e o trabalho em equipe é condição
para a entrega dos projetos, que combinam competências de auditoria, análise
de dados e tecnologia.

\subsection{Importância das soft skills para o desempenho}

Embora a empresa atue em um setor essencialmente técnico, seu sucesso depende
fortemente de \textit{soft skills}. A própria atividade-fim, que é
conscientizar pessoas sobre segurança, exige habilidades de comunicação,
didática e empatia para traduzir conceitos técnicos a públicos não
especializados. Destacam-se como competências comportamentais críticas:

\begin{itemize}[noitemsep]
  \item Comunicação e didática: essenciais na condução de treinamentos e na
        sensibilização dos colaboradores dos clientes;
  \item Empatia: necessária para abordar erros humanos sem gerar
        constrangimento, criando uma cultura de segurança positiva;
  \item Pensamento crítico e ético: indispensável diante do acesso a
        informações sensíveis;
  \item Colaboração: integra as diferentes especialidades dos sócios em
        soluções coesas.
\end{itemize}

A combinação entre competências técnicas (\textit{hard skills}) e
comportamentais (\textit{soft skills}) é, portanto, o principal diferencial
competitivo da Sentinela Cyber.

% ==================================================================
% 5. EVOLUÇÃO DO PENSAMENTO ADMINISTRATIVO
% ==================================================================
\section{Evolução do Pensamento Administrativo}

\subsection{Relação entre práticas observadas e teorias administrativas}

As práticas da Sentinela Cyber podem ser interpretadas à luz das principais
escolas do pensamento administrativo:

\begin{itemize}[noitemsep]
  \item Administração Científica (Taylor): manifesta-se na padronização de
        metodologias, na busca por eficiência na execução dos projetos e na
        medição objetiva de resultados (indicadores e métricas), herança da
        racionalização do trabalho.
  \item Teoria Clássica (Fayol): fundamenta o uso das funções administrativas
        (planejar, organizar, dirigir e controlar) e da divisão de
        responsabilidades entre as diretorias, ainda que de forma flexível.
  \item Teoria das Relações Humanas (Mayo): é fortemente presente, pois a
        empresa valoriza a motivação, o clima organizacional, a comunicação e os
        aspectos humanos, alinhada à própria essência de um negócio centrado no
        fator humano.
  \item Escola Sistêmica: reflete-se na visão da segurança como um sistema
        integrado de pessoas, processos e tecnologia, e na compreensão de que a
        empresa interage continuamente com um ambiente externo dinâmico (novas
        ameaças, regulações e clientes).
\end{itemize}

\subsection{Conceitos ainda aplicados e conceitos superados}

Entre os conceitos ainda aplicados, destacam-se a padronização de processos e a
mensuração de resultados (de inspiração taylorista), as funções administrativas
de Fayol e, sobretudo, a valorização do fator humano proposta pela Escola de
Relações Humanas, plenamente aderente ao modelo de negócio.

Entre os conceitos superados, observa-se o afastamento da rigidez hierárquica e
da visão mecanicista do trabalhador (\textit{homo economicus}) características da
abordagem clássica. A estrutura horizontalizada, a autonomia e a liderança
participativa adotadas pela empresa contrastam com o comando centralizado e a
especialização extrema preconizados pela administração científica em sua forma
original.

% ==================================================================
% 6. REFLEXÃO CRÍTICA
% ==================================================================
\section{Reflexão Crítica}

\subsection{Teoria estudada versus prática observada}

A análise revela que a Sentinela Cyber representa um híbrido entre teoria e
prática. Por um lado, sustenta-se sobre princípios clássicos consolidados, como
planejamento, organização, direção e controle, demonstrando que tais funções
permanecem relevantes mesmo em uma empresa jovem e tecnológica. Por outro, sua
dinâmica cotidiana aproxima-se muito mais das abordagens humanística e
sistêmica, evidenciando que a gestão contemporânea valoriza pessoas,
flexibilidade e adaptação ao ambiente.

Observa-se, contudo, uma lacuna entre a teoria e a prática no campo da
formalização: muitos processos ainda dependem do conhecimento tácito dos
sócios, o que contraria os princípios de documentação e padronização defendidos
pela teoria administrativa e representa um risco à escalabilidade.

\subsection{Propostas de melhoria e recomendações de gestão}

Com base no diagnóstico, recomendam-se as seguintes ações de gestão:

\begin{enumerate}[label=\arabic*.]
  \item Formalizar e documentar processos, reduzindo a dependência das
        pessoas-chave e viabilizando o crescimento sustentável;
  \item Planejar a expansão da equipe, contratando analistas para aliviar a
        sobrecarga dos sócios e permitir foco estratégico;
  \item Implantar um planejamento estratégico formal com indicadores (KPIs)
        claros para a própria empresa, e não apenas para os clientes;
  \item Estruturar um programa de capacitação contínua, em coerência com o pilar
        de Evolução Contínua e com a velocidade das ameaças;
  \item Desenvolver métricas de valor percebido pelos clientes, de modo a
        demonstrar objetivamente o retorno dos serviços comportamentais.
\end{enumerate}

% ==================================================================
% CONCLUSÃO
% ==================================================================
\section{Conclusão}

O estudo da Sentinela Cyber Consultoria permitiu articular os conceitos da
disciplina de Administração de Negócios com a realidade de uma empresa de
consultoria em segurança da informação. Verificou-se que, embora o negócio seja
recente e intensivo em tecnologia, ele se apoia em funções administrativas
clássicas e, simultaneamente, incorpora de maneira intensa as contribuições das
escolas de Relações Humanas e Sistêmica, coerentes com um modelo de negócio
centrado nas pessoas.

A complementaridade dos fundadores, nas áreas de auditoria, análise de dados e
tecnologia, e a cultura organizacional ancorada em quatro pilares (Defesa
Coletiva, Confiança em Rede, Integridade Inegociável e Evolução Contínua)
constituem o principal diferencial competitivo da empresa. Os desafios
identificados, em especial a dependência das pessoas-chave e a baixa
formalização de processos, apontam para um caminho de amadurecimento gerencial.

Conclui-se que a evolução do pensamento administrativo não é meramente
histórica, mas cumulativa: as boas práticas de cada escola coexistem e se
complementam na gestão contemporânea. Para a Sentinela Cyber, o desafio futuro
consiste em equilibrar a flexibilidade e o capital humano que a tornam
competitiva com a estruturação de processos que assegurem sua escalabilidade e
perenidade.

\end{document}
